\documentclass[12pt]{article}
\usepackage[utf8]{inputenc}
\usepackage{eso-pic,graphicx}
\usepackage[a4paper,left=2cm,right=2cm,top=2cm,bottom=2cm, footnotesep=3cm]{geometry}
\usepackage{titlesec}
\usepackage{hyperref}
\usepackage{amsmath}
\usepackage{amsthm}
\usepackage{float}
\usepackage{subcaption}
\usepackage[pages=some]{background}

\newcommand{\HRule}{\rule{\linewidth}{0.3mm}}


% -----------------------------------------------------

\begin{document}

\clearpage

%Title
\textsc{\huge Security Analysis}\\[1cm]

\HRule \\[0.7cm]

\HRule \\[1.2cm]
\vspace{1cm}


\vfill



\clearpage


\tableofcontents
\newpage

% -----------------------------------------------------

\section{Introduction}

This project aims to create a service that analyzes the security of diverse
elements. The implementation is divided in many phases, using iterative development
in order to learn and implement new features and technologies step by step.

The final goal would be to create a full-stack security analysis service, capable of
giving complete reports of the security of a given system, including but not limited
to: web applications, network services, and software applications. The target audience
would be entreprises, small startups and individuals that want to analyze the security of their systems and applications.

Through this project, I aim to develop my cybersecurity skills, as well as my
overall programming skills in full-stack. I'm also planning to deploy the service
in a cloud environment, thus learning about cloud computing and deployment.
Additionally, in much later phases, I would also like to add monetization features
and build a small startup/entreprise around it. This last part is still in the optic
of learning in the sense that many similar services probably exist and my chances
of making concrete profit are really low.

To conclude this introduction, I hope this project will see the light of day as a
real service. I will do my best to document each phase, not only for me but also for
anyone that might read this in the future and is curious about the insights of the
building process of this service (instead of only having the classical README file).

\section{Phase 1 -- Basic Website vulnerability scanner}


\end{document}
